\documentclass[11pt,a4paper,oneside,openany]{book}
\usepackage{biblioflow-ista}

\renewcommand{\titolodocumento}{Post-Modification Test Document}
\renewcommand{\sottotitolodocumento}{Post-Modification Test Document}

\begin{document}
\frontmatter
\copertina
\tableofcontents
\listoftables

\mainmatter
\chapter{Regression Testing}

\section{Regression Testing per la Change Request 01 (CR-01)}

Durante l'implementazione della CR-01 è stata eseguita un'analisi per
identificare le regressioni introdotte dalle modifiche alle route di
prenotazione e per distinguerle dalle differenze volute.

\subsection{Metodologia Operativa}

Il processo di regression testing è stato condotto seguendo un approccio
iterativo per garantire che le nuove funzionalità non compromettessero la
stabilità del sistema esistente. Il flusso di lavoro seguito è stato il
seguente:
\begin{enumerate}
  \item \textbf{Implementazione delle Modifiche:} sono state applicate le
        modifiche al codice di produzione necessarie per la CR-01.
  \item \textbf{Esecuzione della Suite di Test:} è stata rieseguita l'intera
        suite di caratterizzazione della versione di partenza (34 casi).
  \item \textbf{Identificazione dei Fallimenti:} l'esecuzione ha evidenziato
        fallimenti nei casi che descrivevano le route di prenotazione.
  \item \textbf{Analisi delle Cause:} per ogni fallimento è stato stabilito se
        derivasse da una regressione o da un cambiamento richiesto dalla CR.
  \item \textbf{Classificazione:} ogni caso è stato classificato come
        \emph{invariato}, \emph{modificato} o \emph{eliminato}; i casi
        modificati conservano il proprio identificativo e indicano il test che
        ne verifica il nuovo comportamento.
  \item \textbf{Verifica Finale:} i casi invariati sono stati rieseguiti e
        confrontati uno per uno con l'esito della versione di partenza.
\end{enumerate}

\subsection{Caso 1: Regressione nei Test delle Route di Prenotazione}

\textbf{File:} \file{tests/pre-modifica/prenotazioni.test.ts}

Questi casi (TC-PRE-001\dots010) invocano direttamente i route handler di
prenotazione sostituendo Prisma e la sessione con mock costruiti sulla
versione di partenza.

\paragraph{Il Problema}
Dopo la CR i casi falliscono. Le route ottengono l'utente con
\texttt{requireUser()} e delegano la creazione a un servizio che apre una
transazione: i mock storici non esportano né le nuove funzioni di
autenticazione né \texttt{\$transaction}, quindi l'esecuzione si interrompe
prima delle asserzioni.

\paragraph{Soluzione Implementata}
I mock storici non sono stati riscritti: avrebbero cancellato l'evidenza del
comportamento originale, compresi i tre difetti fissati da TC-PRE-004,
TC-PRE-007 e TC-PRE-010. I casi sono classificati come \textbf{modificati} e
ciascuno rimanda alla suite che verifica il comportamento nuovo; per i tre
difetti la differenza è voluta.

{\small
\begin{xltabular}{\linewidth}{L{2.1cm}YL{4.3cm}}
  \toprule
  \intestazione \tb{Caso} & \tb{Nuovo comportamento} & \tb{Test che lo verifica} \\
  \midrule
  \endhead
  TC-PRE-001 & identità dalla sessione, creazione in transazione
  \texttt{Serializable} & \texttt{TC-BIB31-001}, auth-prenotazioni \\
  TC-PRE-002, 003 & validazione centralizzata prima di ogni scrittura &
  \texttt{TC-BIB27-003}\dots\texttt{013} \\
  TC-PRE-004 & intervallo invertito rifiutato con 422 (prima 201) &
  \texttt{TC-BIB27-005} \\
  TC-PRE-005 & conflitto garantito dal vincolo del database, con proposta di
  coda & \texttt{TC-BIB24-DB-001}\dots\texttt{003}, \texttt{TC-BIB31-002} \\
  TC-PRE-006 & cancellazione solo del titolare, con promozione della coda &
  \texttt{TC-CODA-DEL-001}\dots\texttt{005} \\
  TC-PRE-007, 008 & check-in con data e orario ricomposti, sessione e
  titolarità su entrambi i percorsi & finestra di check-in
  (\texttt{TC-FIN-*}), auth-prenotazioni \\
  TC-PRE-009, 010 & estensione con le regole centralizzate; durata calcolata
  sul giorno della prenotazione & \texttt{TC-BIB27-006}\dots\texttt{008},
  \texttt{015}\dots\texttt{019} \\
  \bottomrule
  \caption{Casi di prenotazione modificati dalla CR-01}\label{tab:modificati-cr01}\\
\end{xltabular}}

\subsection{Caso 2: Regressione nel Contratto della Pagina di Prenotazione}

\textbf{File:} \file{tests/pre-modifica/admin-ui.test.ts} (TC-PRE-023,
TC-PRE-024)

\paragraph{Il Problema}
TC-PRE-023 verifica che la pagina \texttt{/prenota} invii
\texttt{userId}, posto, data e orari. Dopo la CR il caso continua a passare:
la pagina invia ancora l'identificativo. Il rischio era interpretare questo
esito come conferma che il server se ne fidi ancora.

\paragraph{Soluzione Implementata}
Il caso è classificato come \textbf{modificato}: il contratto della pagina è
invariato, ma il server ignora l'identificativo ricevuto e usa quello della
sessione. Lo dimostra il test di integrazione \enquote{ignora userId del
payload e crea per l'utente autenticato}, che invia l'identificativo di un
altro studente e verifica che la prenotazione sia intestata a chi ha la
sessione. TC-PRE-024 rimanda alla stessa prova e ai casi di ingresso in coda
dal conflitto (\texttt{TC-BIB52}\dots\texttt{54}).

\subsection{Analisi di Impatto su Accesso e Area Staff}

\textbf{File:} \file{tests/pre-modifica/login.test.ts},
\file{tests/pre-modifica/admin-ui.test.ts}

I casi di accesso, statistiche, gestione utenti e prestiti (TC-PRE-011\dots020
e TC-PRE-022) non hanno subito regressioni. Le nuove funzioni di
autorizzazione sono state aggiunte accanto a quelle esistenti, senza
cambiarne il comportamento: il provider di accesso e le guardie di ruolo
dell'area staff sono rimasti compatibili con i casi storici.

\section{Regression Testing per la Change Request 02 (CR-02)}

La CR-02 aggiunge un'entità e nuovi percorsi (coda, promozione, scadenza), ma
interviene anche su punti esistenti: l'annullamento dello staff, il rilascio
automatico dei no-show e il canale in tempo reale. Il regression testing ha
verificato che queste estensioni non alterassero i comportamenti preesistenti.

\subsection{Strategia di Regressione}

La strategia si è concentrata su due momenti:
\begin{enumerate}
  \item \textbf{Chiusura delle Change Requests:} riesecuzione dei 20 casi
        classificati come invariati, con confronto caso per caso con l'esito
        della versione di partenza.
  \item \textbf{Versione finale:} riesecuzione degli stessi casi dopo la
        manutenzione correttiva successiva alle Change Requests, per
        distinguere le regressioni dalle correzioni intenzionali.
\end{enumerate}

\subsection{Esecuzione della Suite alla Chiusura delle Change Requests}

I 20 casi invariati sono stati rieseguiti sulla versione che chiude
l'implementazione (commit \texttt{adc5e30}, 4 settembre 2026) e confrontati con
l'esecuzione della versione di partenza.

\textbf{Risultati:}
\begin{itemize}
  \item \textbf{Casi invariati:} 20 passati su 20, identici alla versione di
        partenza.
  \item \textbf{Aree verificate:} accesso (TC-PRE-011\dots013), protezione di
        pagine e API (TC-PRE-014\dots016), area staff (TC-PRE-017\dots020,
        TC-PRE-022), automazioni e notifiche preesistenti (PRE-AUT-001, 002,
        004, PRE-NOT-001, 002), contratto del canale SSE (PRE-SSE-001\dots004).
\end{itemize}

Due casi modificati meritano una nota: TC-PRE-021 (annullamento dello staff)
e PRE-AUT-003 (no-show automatico) continuano a passare anche dopo la CR,
perché con la coda vuota la promozione non produce effetti. La scelta di
attivare la promozione solo se esiste qualcuno in attesa ha garantito la
compatibilità con il comportamento precedente.

\subsection{Analisi sulla Versione Finale}

Dopo le Change Requests il progetto ha attraversato una fase di manutenzione
correttiva. La riesecuzione degli stessi 20 casi sulla versione finale dà il
seguente quadro:

{\small
\begin{xltabular}{\linewidth}{L{3.3cm}L{2.2cm}Y}
  \toprule
  \intestazione \tb{Casi} & \tb{Classificazione} & \tb{Motivo} \\
  \midrule
  \endhead
  TC-PRE-011\dots015, 017\dots019, 022, PRE-AUT-004, PRE-NOT-001, 002 &
  Invariato, \ok & -- \\
  TC-PRE-016 & Modificato &
  Il middleware ora verifica davvero il token: un cookie inventato non apre
  più \texttt{/prenota} (si veda il Test Incident n.4). \\
  TC-PRE-020 & Modificato &
  La restituzione registrata dallo staff ora incrementa anche le copie
  disponibili del libro in una transazione; prima ogni restituzione al banco
  \enquote{perdeva} una copia. Il mock storico non prevede la transazione;
  il comportamento nuovo è verificato su database reale. \\
  PRE-AUT-001, 002 & Modificato &
  Gli orari delle prenotazioni sono ora salvati nell'ora di Roma, mentre la
  fixture storica li scrive in UTC: il promemoria non scatta più all'orario
  della fixture. Il comportamento è verificato da
  \file{tests/unit/automation-service.test.ts}. \\
  PRE-SSE-001\dots004 & Eliminato &
  La catena SSE non è mai stata collegata all'interfaccia e, tenendo i client
  in memoria, non funziona con più istanze serverless. È stata rimossa e
  sostituita dalla rilettura autenticata delle API e dal controllo periodico
  delle notifiche. \\
  \bottomrule
  \caption{Casi invariati alla chiusura delle CR, rieseguiti sulla versione finale}\label{tab:regressione-finale}\\
\end{xltabular}}

\subsection{Conclusione}

Il Regression Testing ha dato esito positivo in entrambi i momenti. Alla
chiusura delle Change Requests tutti i 20 casi invariati conservano l'esito
della versione di partenza. Sulla versione finale restano 12 casi invariati,
tutti superati; gli altri 8 sono stati modificati o eliminati da interventi
documentati, ciascuno con un test che ne verifica il nuovo comportamento.

\chapter{Copertura}

In questo capitolo vengono presentati i risultati della copertura del codice
(code coverage) per il codice server e per l'interfaccia di BiblioFlow.

\section{Metodologia}

Per garantire una valutazione completa dell'efficacia della suite di test,
sono state misurate sia la Statement Coverage (copertura delle istruzioni) sia
la Branch Coverage (copertura dei rami).

\subsection{Copertura Backend}

La copertura del codice server (route handler, librerie, middleware, hook) è
stata calcolata con Vitest e il fornitore di copertura nativo di V8, che
misura anche i rami per individuare i percorsi sfuggiti a una semplice
verifica delle righe.
\begin{itemize}
  \item \textbf{Strumento:} Vitest 4.1.11 con \texttt{@vitest/coverage-v8}.
  \item \textbf{Comando:}
\begin{lstlisting}
npm run test:coverage
\end{lstlisting}
  \item \textbf{Configurazione:} la misura include tutto \file{src}, route
        handler compresi, ed esclude la suite di caratterizzazione. Sono
        configurate soglie minime: 90\% di istruzioni e 80\% di rami su
        \file{src/lib} e \file{src/app/api}, 100\% sul middleware; sotto
        soglia il comando fallisce.
\end{itemize}

\subsection{Copertura Frontend}

I componenti React toccati dalle Change Requests sono verificati in un DOM
simulato (jsdom) con Testing Library, nella stessa esecuzione di Vitest.
\begin{itemize}
  \item \textbf{Strumento:} Vitest, jsdom 26, Testing Library.
  \item \textbf{Configurazione:} test dei componenti in
        \file{tests/unit/componenti}. Le pagine Next.js sono esercitate dalla
        suite E2E su build di produzione, che non contribuisce alla misura V8.
\end{itemize}

\subsection{Risultati}

La tabella seguente riassume le metriche di copertura ottenute, confrontate
con quelle misurate all'inizio dell'attività di miglioramento della suite.

{\small
\begin{xltabular}{\linewidth}{YC{1.9cm}C{1.9cm}C{2.9cm}}
  \toprule
  \intestazione \tb{Componente} & \tb{Statement} & \tb{Branch} & \tb{Prima (stmt / branch)} \\
  \midrule
  \endhead
  Route handler (\file{src/app/api}, 37 route) & \textbf{95,1\%} & \textbf{89,0\%} & 72,3\% / 66,6\% \\
  Librerie (\file{src/lib}, 29 moduli) & \textbf{95,4\%} & \textbf{90,5\%} & 89,9\% / 83,6\% \\
  Middleware & 100\% & 100\% & 100\% / 100\% \\
  Hook di ingresso in coda & 94,1\% & 84,4\% & 25,5\% / 25,0\% \\
  \textbf{Totale codice server} & \textbf{95,2\%} & \textbf{89,4\%} & -- \\
  Moduli toccati dalle CR (18 file, UI compresa) & 87,4\% & 80,3\% & -- \\
  Componenti (\file{src/components}) & 21,4\% & 17,9\% & circa 3\% \\
  \bottomrule
  \caption{Metriche di Copertura del Codice}\label{tab:copertura}\\
\end{xltabular}}

La suite conta 123 file e 1013 test, tutti superati: 815 test di unità, 32
test di componenti, 138 test di integrazione e 28 test sui vincoli del
database. Prima dell'attività i test erano 763 in 103 file.

\section{Analisi}

\textbf{Codice server:} il codice server ha raggiunto un grado di copertura
elevato (95,2\% delle istruzioni, 89,4\% dei rami), sopra le soglie del Master
Test Plan. È importante notare che la configurazione iniziale escludeva dalla
misura proprio i route handler: il dato pubblicato in precedenza non
comprendeva le API. Misurate per la prima volta, le route si fermavano al
72,3\% delle istruzioni, con sette route mai eseguite da alcun test
(registrazione, verifica e reinvio dell'email, profilo, dettaglio di un posto
e di un libro, storico di un utente).

Di conseguenza, sono stati introdotti test mirati per colmare questo
divario:
\begin{enumerate}
  \item \textbf{Route a copertura zero o bassa:} test per registrazione e
        verifica dell'email su database reale, per dettaglio posti e utenti,
        anomalie, prestiti dello staff, catalogo, sale, profilo e
        notifiche, inclusi accessi senza sessione e con ruolo errato.
  \item \textbf{Errori interni:} un test trasversale chiama 49 handler con il
        database guasto e verifica che rispondano 500 con un messaggio
        generico, senza nomi di tabelle, driver o stack.
  \item \textbf{Ciclo di vita della coda e modifica dello staff} su database
        reale, componenti dell'interfaccia in jsdom e una suite E2E.
\end{enumerate}

I nuovi test hanno fatto emergere quindici difetti, tutti corretti e coperti
da un test che fallisce sul codice precedente:

{\small
\begin{xltabular}{\linewidth}{L{1.5cm}YL{3.1cm}}
  \toprule
  \intestazione \tb{ID} & \tb{Difetto} & \tb{Origine} \\
  \midrule
  \endhead
  D-01 & Registrazione ed estensione registrate nel log come
  \enquote{prenotazione creata}. & Baseline; estensione dalla CR-01 \\
  D-02 & Nel dettaglio di un posto, \enquote{oggi} calcolato nel fuso del
  server: tra mezzanotte e le due di Roma la prenotazione in corso spariva. &
  Baseline \\
  D-03 & I no-show rilevati dal cron non comparivano tra le anomalie né nei
  contatori dello staff. & Baseline \\
  D-04 & \enquote{Risolvi tutte} chiudeva un solo evento per utente. & Baseline \\
  D-05 & Solleciti e rinnovi registrati come restituzioni e nuovi prestiti; il
  sollecito singolo non verificava lo stato del prestito. & Baseline \\
  D-06 & Durata del prestito decisa dal client, anche negativa o di anni. & Baseline \\
  D-07\dots09 & Filtri non validi (stato del prestito, tipo di notifica, date
  delle prenotazioni) producevano un errore 500 invece di 422. & Baseline \\
  D-10 & Una modifica dello staff in sovrapposizione, rifiutata dal vincolo
  del database, rispondeva \enquote{errore interno} invece di 409. &
  Effetto a catena della CR-01 \\
  D-11 & \enquote{1 disponibili}, \enquote{1 risultati disponibili}. & Baseline \\
  D-12 & Pulsanti dello staff fatti solo di icona senza nome accessibile;
  azioni proposte su prenotazioni già concluse. & Baseline e manutenzione
  correttiva \\
  D-13 & Nessuna schermata mostrava allo studente le proprie richieste in
  coda né permetteva di uscirne. & CR-02 (interfaccia incompleta) \\
  D-14 & Le prenotazioni scadute per promozione non confermata non
  comparivano nello storico. & CR-02 \\
  D-15 & Una data non valida nell'indirizzo della pagina staff delle
  prenotazioni causava un errore del server. & Baseline \\
  \bottomrule
  \caption{Difetti emersi dall'aumento della copertura}\label{tab:difetti}\\
\end{xltabular}}

\textbf{Interfaccia:} la copertura dei componenti passa da circa il 3\% al
21,4\%. Il valore riguarda tutti i componenti dell'applicazione: quelli
toccati dalle Change Requests (mappa, griglia mobile, lista d'attesa, azioni
dello staff) sono verificati in dettaglio, mentre le pagine sono esercitate
dalla suite E2E. Per questo motivo il totale grezzo dell'intero \file{src}
(48,5\%) non è indicativo e non viene usato come soglia.

\chapter{Unit Testing}

\section{Unit Testing per la Change Request 01 (CR-01)}

Questa sezione descrive i test di unità introdotti per la CR-01: prenotazione
sicura e atomica. L'obiettivo era verificare in isolamento le regole del
servizio di dominio, le funzioni di autorizzazione e il comportamento delle
route.

\subsection{Tecnologie e Framework}

Per lo sviluppo della suite di test sono stati impiegati i seguenti strumenti:
\begin{itemize}
  \item \textbf{Vitest:} framework di testing primario, per la definizione dei
        casi, le asserzioni e la misura della copertura.
  \item \textbf{Mock di Vitest} (\texttt{vi.mock}, \texttt{vi.hoisted}):
        sostituiscono il client Prisma e la sessione di Auth.js, isolando la
        logica dal database e dall'autenticazione.
  \item \textbf{Richieste HTTP native:} le route vengono invocate con oggetti
        \texttt{Request} costruiti nel test, esattamente come le chiamerebbe
        Next.js.
\end{itemize}

\subsection{Logica Core: servizio di dominio delle prenotazioni}

\textbf{File:} \file{tests/unit/prenotazioni-service.test.ts} (38 casi)

\begin{itemize}
  \item \texttt{TC-BIB27-001}\dots\texttt{020}: validazione centralizzata
        degli intervalli, tra cui fine non successiva all'inizio
        (\texttt{005}), durata minima di un'ora e massima di otto
        (\texttt{006}\dots\texttt{008}), orari della sala e sala inattiva
        (\texttt{012}, \texttt{013}), conflitti sul posto e sull'utente con
        proposta di coda (\texttt{015}\dots\texttt{019}).
  \item \texttt{TC-BIB31-001}: la creazione avviene in una transazione
        \texttt{Serializable}.
  \item \texttt{TC-BIB31-002}: la violazione del vincolo del database viene
        tradotta in un conflitto proponibile come coda, non in un errore
        interno.
\end{itemize}

\subsection{Funzioni di Autorizzazione}

\textbf{File:} \file{tests/unit/auth-helpers.test.ts}

\begin{itemize}
  \item \texttt{TC-BIB26-001}\dots\texttt{007}: \texttt{requireUser},
        \texttt{requireRole} e \texttt{assertOwnership} restituiscono 401
        senza sessione, 403 per il ruolo non ammesso e 404 per la risorsa di
        un altro utente; lo staff è sempre ammesso.
\end{itemize}

\subsection{Route di Prenotazione}

\textbf{File:} \file{tests/unit/prenotazioni-elenco-e-creazione.test.ts},
\file{tests/unit/prenotazioni-regole-finestra-checkin.test.ts},
\file{tests/unit/errori-generici-route.test.ts}

\begin{itemize}
  \item L'elenco è sempre ristretto alla sessione, anche se l'indirizzo
        contiene l'identificativo di un altro utente; filtri di data non
        validi producono 422.
  \item La creazione ignora l'identificativo inviato dal client e, in caso
        di conflitto, risponde 409 con la proposta di entrare in coda.
  \item La finestra di check-in è verificata ai suoi estremi: quindici minuti
        prima e dopo l'inizio, trenta per chi ha il margine pendolare
        (\texttt{TC-FIN-*}, \texttt{TC-TOL-*}).
  \item Con il database guasto ogni route risponde 500 senza dettagli
        interni.
\end{itemize}

\section{Unit Testing per la Change Request 02 (CR-02)}

La CR-02 introduce la lista d'attesa e la promozione automatica. La strategia
di testing unitario ha coperto tutti i livelli coinvolti: servizi,
automazioni, route e interfaccia.

\subsection{Stack Tecnologico di Testing}

Data la natura dell'applicazione, sono stati adottati strumenti specifici per
ogni livello:
\begin{itemize}
  \item \textbf{Vitest con mock:} per i servizi, le automazioni e le route,
        isolando database, orologio e sessione.
  \item \textbf{Testing Library, user-event e jsdom:} per i componenti React.
        Permettono di verificare il rendering, lo stato e le interazioni in un
        DOM simulato, interrogando gli elementi per ruolo e nome accessibile
        come farebbe una tecnologia assistiva.
\end{itemize}

\subsection{Servizi e Automazioni}

\textbf{File:} \file{tests/unit/automation-service.test.ts},
\file{tests/unit/cron-automations.test.ts},
\file{tests/unit/prenotazioni-cancellazione-promuove-coda.test.ts},
\file{tests/unit/admin-prenotazioni.test.ts}

\begin{itemize}
  \item \texttt{TC-BIB40}\dots\texttt{46}: il no-show avvia la promozione;
        le notifiche di ingresso, uscita, promozione e scadenza; la scadenza
        delle promozioni non confermate dopo quindici minuti; la correlazione
        tra notifiche ed eventi.
  \item \texttt{TC-BIB41-001}\dots\texttt{003}: il lock impedisce due
        esecuzioni contemporanee del cron.
  \item \texttt{TC-CODA-DEL-001}\dots\texttt{005} e
        \texttt{TC-CODA-PATCH-*}: la cancellazione dello studente promuove il
        primo in coda per lo stesso intervallo, solo dopo che la prenotazione
        risulta annullata, e notifica il promosso.
  \item \texttt{TC-BIB49-001}\dots\texttt{005}: l'annullamento dello staff
        promuove e restituisce il nome del promosso all'interfaccia.
\end{itemize}

\subsection{Frontend Testing (React)}

\textbf{File:} \file{tests/unit/componenti/*.test.tsx},
\file{tests/unit/mappa-coda.test.ts}, \file{tests/unit/mobile-coda.test.ts},
\file{tests/unit/prenota-conflitto.test.ts}

\begin{itemize}
  \item \textbf{Mappa dei posti:} legenda, posti interattivi, uso da tastiera
        (Invio e spazio), ingresso in coda con lo stesso contratto della vista
        mobile, ripiego sulla griglia sotto 768 pixel.
  \item \textbf{Griglia mobile:} ordine naturale dei posti, nome accessibile e
        stato premuto, posti in manutenzione non selezionabili, ingresso in
        coda, rifiuto del server, assenza di rete.
  \item \textbf{Lista d'attesa dello studente:} elenco con posizione,
        uscita dalla coda con rilettura delle posizioni, stato vuoto, errore
        con \enquote{Riprova} (\texttt{TC-UI-CODA-001}\dots\texttt{005}).
  \item \textbf{Azioni dello staff:} azioni proposte solo per gli stati
        accettati dal server, esito della promozione dopo l'annullamento,
        conflitto 409 mostrato all'operatore.
  \item \textbf{Conflitto all'ultimo istante:} il messaggio di posto non più
        disponibile propone la lista d'attesa (\texttt{TC-BIB54-001}\dots\texttt{004}).
\end{itemize}

\subsection{Middleware della Coda}

\textbf{File:} \file{tests/unit/middleware-coda.test.ts}

\begin{itemize}
  \item \texttt{TC-BIB51-001}\dots\texttt{005}: la route della coda è
        dichiarata nel matcher e respinge le richieste senza sessione (401)
        o con un cookie inventato; un token firmato davvero viene accettato e
        la protezione delle API dello staff resta invariata.
\end{itemize}

\chapter{Integration Testing}

\section{Integration Testing per la Change Request 01 (CR-01)}

Questa sezione documenta le attività di integration testing svolte per
validare la CR-01. Mentre lo unit testing isola i singoli componenti,
l'integration testing verifica che route, servizio di dominio e database
collaborino correttamente e producano l'esito atteso in uno scenario reale.

\subsection{Strategia di Testing}

Per i test di integrazione è stata adottata una strategia black-box parziale:
\begin{itemize}
  \item \textbf{Ambiente Reale:} i test usano PostgreSQL 16 reale, con le
        migrazioni applicate e un database isolato il cui nome deve
        dichiarare di essere di test; i dati vengono ripuliti prima di ogni
        caso.
  \item \textbf{Esecuzione dal Route Handler:} le richieste attraversano la
        route, il servizio di dominio e il database, senza mock di Prisma.
  \item \textbf{Validazione dell'Esito:} si verificano la risposta HTTP e lo
        stato finale del database (prenotazioni, eventi, notifiche).
\end{itemize}

\subsection{Suite di Test}

\paragraph{Concorrenza} (\file{tests/integration/concorrenza.test.ts})
\begin{itemize}
  \item \textbf{Due richieste simultanee, tre round:} sullo stesso posto e
        intervallo esattamente una richiesta ottiene la prenotazione.
  \item \textbf{N richieste simultanee:} un solo vincitore e nessun record
        fantasma lasciato dalle transazioni rifiutate.
\end{itemize}

\paragraph{Vincolo del database}
(\file{tests/post-modifica/prenotazioni-constraint-db.test.ts})
\begin{itemize}
  \item \texttt{TC-BIB24-DB-001}\dots\texttt{003}: il vincolo di esclusione,
        verificato scrivendo direttamente sul database, considera occupato lo
        slot anche durante il check-in, ammette intervalli adiacenti e ignora
        le prenotazioni in stato terminale.
\end{itemize}

\paragraph{Autenticazione e titolarità}
(\file{tests/integration/auth-prenotazioni.test.ts},
\file{tests/integration/autorizzazione-admin.test.ts})
\begin{itemize}
  \item 401 su ogni operazione protetta senza sessione; 404 su lettura,
        modifica e cancellazione della prenotazione di un altro studente.
  \item Identificativo del payload ignorato: la prenotazione è intestata a
        chi ha la sessione.
  \item \texttt{TC-BIB57-*} (36 casi): 401 e 403 su tutte le route dello
        staff, comprese le tre rimaste aperte alla prima verifica.
\end{itemize}

\section{Integration Testing per la Change Request 02 (CR-02)}

L'integration testing per la CR-02 ha l'obiettivo di verificare la catena che
parte da un evento di liberazione del posto e arriva alla promozione, alle
notifiche e agli eventi registrati.

\subsection{Strumenti Utilizzati}

\begin{itemize}
  \item \textbf{Vitest:} runner dell'intera suite.
  \item \textbf{PostgreSQL 16:} database reale effimero, in locale e nel job
        \texttt{test} della pipeline.
  \item \textbf{Prisma Client:} usato dai test per preparare i dati e
        verificare lo stato finale.
\end{itemize}

\subsection{Componenti Sotto Test}

I test di integrazione coinvolgono la cooperazione tra:
\begin{itemize}
  \item \textbf{Route handler:} cancellazione dello studente e dello staff,
        coda, cron.
  \item \textbf{Servizio di dominio:} ingresso, uscita, posizione e
        promozione.
  \item \textbf{Servizio delle automazioni:} no-show, scadenza delle
        promozioni, lock.
  \item \textbf{Database:} vincolo di unicità delle richieste attive e
        ordinamento della coda.
\end{itemize}

\subsection{Suite di Test}

\paragraph{Test 1: Ciclo coda--promozione}
(\file{tests/integration/coda-promozione.test.ts})
\begin{enumerate}
  \item B entra in coda per il posto prenotato da A.
  \item Lo staff annulla la prenotazione di A.
  \item B viene promosso con prenotazione, notifica ed evento correlati
        (\texttt{TC-BIB58-001}); con la coda vuota nessuno viene promosso
        (\texttt{TC-BIB58-002}).
\end{enumerate}

\paragraph{Test 2: Automazioni} (\file{tests/integration/automazioni.test.ts})
\begin{enumerate}
  \item Il no-show automatico promuove il primo in coda (\texttt{TC-BIB47-001}).
  \item Due esecuzioni ravvicinate non producono promozioni, notifiche o
        richieste duplicate (\texttt{002}, \texttt{002b}); due invocazioni
        concorrenti del cron eseguono una sola volta (\texttt{003}).
  \item La promozione non confermata scade e il posto passa al successivo
        (\texttt{005}); la scadenza è idempotente (\texttt{006}) e non tocca
        una promozione confermata con il check-in (\texttt{007}).
\end{enumerate}

\paragraph{Test 3: Ciclo di vita della coda}
(\file{tests/integration/coda-ciclo-di-vita-db.test.ts},
\file{tests/post-modifica/lista-attesa-db.test.ts})
\begin{enumerate}
  \item Tre utenti ottengono le posizioni 1, 2 e 3 in ordine di arrivo;
        non si entra due volte nella stessa coda.
  \item Il titolare esce dalla coda: resta traccia nel log e chi era dietro
        avanza; la richiesta di un altro utente risulta inesistente.
  \item Dopo l'uscita si può rientrare, in fondo; con lo slot ancora occupato
        nessuno viene promosso.
\end{enumerate}

\subsection{Risultati}

Tutti i 166 test di integrazione e sui vincoli del database sono stati
eseguiti con successo, confermando che promozione, notifiche ed eventi
restano coerenti anche con esecuzioni ripetute e concorrenti.

\chapter{System Testing}

A seguito dell'implementazione delle Change Request CR-01 (prenotazione
sicura e atomica) e CR-02 (lista d'attesa e riassegnazione), la strategia di
System Testing è stata aggiornata. Oltre alla riesecuzione dei casi esistenti
descritti nel documento \enquote{Pre-Modification Test Document}, sono stati
introdotti nuovi casi di test per coprire le nuove funzionalità.

\section{Aggiornamento Categorie}

Le categorie sono state estese per includere i parametri rilevanti per la
titolarità, la concorrenza e la lista d'attesa. Oltre alle categorie già
definite nel documento pre-modifica, sono state aggiunte le seguenti.

\subsection{Nuovo Parametro: Titolarità e Concorrenza}

{\small
\begin{xltabular}{\linewidth}{L{5.4cm}L{1.3cm}Y}
  \toprule
  \intestazione \tb{Categoria} & \tb{Codice} & \tb{Scelte (Classi di Equivalenza)} \\
  \midrule
  \endhead
  \multirow{2}{=}{\textbf{Titolarità (TI)}\newline\footnotesize Proprietario della risorsa richiesta.}
    & TI0 & Propria \\*
    & TI1 & Di un altro utente \\
  \midrule
  \multirow{2}{=}{\textbf{Concorrenza (CC)}\newline\footnotesize Richieste contemporanee sullo stesso posto e intervallo.}
    & CC0 & Richiesta singola \\*
    & CC1 & Richieste simultanee \\
  \bottomrule
  \caption{Nuovi parametri: Titolarità e Concorrenza}\label{tab:par-titolarita}\\
\end{xltabular}}

\subsection{Nuovo Parametro: Lista d'Attesa (LA, EL)}

{\small
\begin{xltabular}{\linewidth}{L{5.4cm}L{1.3cm}Y}
  \toprule
  \intestazione \tb{Categoria} & \tb{Codice} & \tb{Scelte (Classi di Equivalenza)} \\
  \midrule
  \endhead
  \multirow{3}{=}{\textbf{Stato della Coda (LA)}\newline\footnotesize Richieste in attesa per il posto e l'intervallo.}
    & LA0 & Coda vuota \\*
    & LA1 & Altri utenti in coda \\*
    & LA2 & Richiesta attiva dello stesso utente \\
  \midrule
  \multirow{5}{=}{\textbf{Evento di Liberazione (EL)}\newline\footnotesize Causa per cui il posto torna libero.}
    & EL0 & Nessuno (slot ancora occupato) \\*
    & EL1 & Cancellazione dello studente \\*
    & EL2 & Annullamento dello staff \\*
    & EL3 & No-show rilevato dal cron \\*
    & EL4 & Promozione non confermata scaduta \\
  \bottomrule
  \caption{Nuovi parametri: Lista d'Attesa}\label{tab:par-coda}\\
\end{xltabular}}

\section{Nuovi Test Case (TC)}

In base alle nuove categorie identificate, sono stati progettati e
implementati i seguenti Test Case aggiuntivi. Ogni caso è automatizzato: la
colonna dell'oracolo indica il test che lo esegue.

{\small
\begin{xltabular}{\linewidth}{L{2.2cm}L{3.4cm}Y}
  \toprule
  \intestazione \tb{TC ID} & \tb{Test Frame (Parametri)} & \tb{Oracolo / Risultato Atteso} \\
  \midrule
  \endhead
  TC-POST-001 & SE0, RS1, OP0\dots OP4 &
  \textbf{Senza sessione:} 401 su ogni operazione sulle prenotazioni e sulla
  coda (auth-prenotazioni, \texttt{TC-BIB57-040}\dots\texttt{042}). \\
  TC-POST-002 & SE1, RU0, TI1, OP1\dots OP4 &
  \textbf{Risorsa altrui:} 404 su lettura, modifica e cancellazione;
  nessuna scrittura (auth-prenotazioni). \\
  TC-POST-003 & SE1, RU0, OP0, OR0, TI1 &
  \textbf{Identità nel payload:} la prenotazione è intestata a chi ha la
  sessione, non all'utente indicato nel corpo (auth-prenotazioni). \\
  TC-POST-004 & SE1, RU0, RS1, FS0\dots FS3 &
  \textbf{Studente sull'area staff:} 403 su tutte le route dello staff
  (\texttt{TC-BIB57-010}\dots\texttt{015}, \texttt{050}\dots\texttt{055}). \\
  TC-POST-005 & SE1, RU0, OP0, IO2 &
  \textbf{Intervallo invertito:} 422, nessuna prenotazione
  (\texttt{TC-BIB27-005}; prima TC-PRE-004 rispondeva 201). \\
  TC-POST-006 & SE1, RU0, OP0, IO0, DI0, CC1 &
  \textbf{Richieste simultanee:} una sola prenotazione creata, le altre 409,
  nessun record fantasma (concorrenza). \\
  TC-POST-007 & SE1, RU0, OP0, IO0, DI1, CC0 &
  \textbf{Posto occupato:} 409 con la proposta di entrare in coda
  (\texttt{TC-BIB31-002}, \texttt{TC-BIB27-015}). \\
  TC-POST-008 & SE1, RU0, LA0, LA1 &
  \textbf{Ingresso in coda:} posizioni 1, 2, 3 in ordine di arrivo
  (\texttt{TC-CODA-VITA-001}). \\
  TC-POST-009 & SE1, RU0, LA2 &
  \textbf{Richiesta duplicata:} rifiutata; resta una sola richiesta attiva
  (\texttt{TC-CODA-VITA-002}, \texttt{TC-BIB23-DB-001}). \\
  TC-POST-010 & SE1, RU0, TI0, LA1 &
  \textbf{Uscita dalla coda:} richiesta \texttt{ANNULLATA}, evento
  registrato, chi era dietro avanza (\texttt{TC-CODA-VITA-010}). \\
  TC-POST-011 & EL1, LA1 &
  \textbf{Cancellazione dello studente:} il primo in coda ottiene una
  prenotazione \texttt{CONFERMATA} e la notifica di promozione
  (\texttt{TC-CODA-DEL-001}\dots\texttt{003}). \\
  TC-POST-012 & SE1, RU1, EL2, LA1 &
  \textbf{Annullamento dello staff:} promozione del primo in coda e nome del
  promosso restituito allo staff (\texttt{TC-BIB58-001}, \texttt{TC-BIB49-*}). \\
  TC-POST-013 & EL2, LA0 &
  \textbf{Coda vuota:} nessuna promozione e nessuna notifica di promozione
  (\texttt{TC-BIB58-002}, \texttt{TC-CODA-DEL-004}). \\
  TC-POST-014 & EL3, LA1, ES0, ES1 &
  \textbf{No-show:} promozione del primo in coda; con due esecuzioni nessuna
  promozione né notifica duplicata (\texttt{TC-BIB47-001}, \texttt{002}). \\
  TC-POST-015 & EL4, LA1 &
  \textbf{Promozione scaduta:} dopo quindici minuti senza check-in la
  promozione scade e il posto passa al successivo (\texttt{TC-BIB47-005}). \\
  TC-POST-016 & EL0, LA1 &
  \textbf{Slot ancora occupato:} nessuno viene promosso
  (\texttt{TC-CODA-VITA-020}). \\
  TC-POST-017 & SE1, RU0, OP2, ST0, MO0 &
  \textbf{Check-in nella finestra:} accettato (prima TC-PRE-007 lo
  rifiutava); estremi della finestra verificati (\texttt{TC-FIN-*}). \\
  TC-POST-018 & SE1, RU0, OP4, ST0, DI0 &
  \textbf{Estensione libera:} accettata se la durata totale non supera otto
  ore (\texttt{TC-BIB27-007}; prima TC-PRE-010 la rifiutava). \\
  \bottomrule
  \caption{Nuovi Test Case di sistema}\label{tab:tc-post}\\
\end{xltabular}}

\chapter{WebApp Testing (End-to-End)}

Parallelamente ai System Test, è stata implementata una suite di test
End-to-End (E2E) automatizzati per validare l'esperienza utente sulla WebApp.
La suite è stata sviluppata con Playwright, che pilota un browser reale e
simula le interazioni dell'utente (click, compilazione dei moduli,
navigazione) in un ambiente controllato.

\section{Strategia di Testing E2E}

I test E2E si concentrano sui flussi delle due Change Requests, verificando
l'integrazione completa tra pagine, API e database sulla build di produzione.
\begin{itemize}
  \item \textbf{Ambiente:} build di produzione di Next.js, Chromium,
        lingua italiana e fuso orario di Roma; PostgreSQL 16 con migrazioni e
        dati dimostrativi caricati prima della suite, su un database il cui
        nome deve dichiarare di essere di test.
  \item \textbf{Esecuzione:} un solo worker, scenari in serie che condividono
        lo stato del database come farebbero utenti reali; tracce e
        screenshot conservati in caso di errore.
  \item \textbf{Accesso e Autorizzazione:} redirect al login e ritorno alla
        pagina richiesta, credenziali errate, ruoli, cookie contraffatti.
  \item \textbf{Happy Path:} prenotazione guidata, coda, promozione.
  \item \textbf{Negative Testing:} posto preso da un altro studente
        all'ultimo istante.
  \item \textbf{Data Binding:} posizione in coda, nome del promosso, notifica
        e prenotazione a nome del promosso.
\end{itemize}

La suite è configurata anche nella pipeline di integrazione continua, in un
job dedicato, ed è risultata stabile su tre esecuzioni consecutive (12 scenari
su 12, circa 37 secondi ciascuna).

\section{Casi di Test E2E Automatizzati}

Di seguito sono riportati i Test Case implementati nei file
\file{tests/e2e/*.spec.ts}. Ogni test è identificato da un ID univoco E2E.

{\small
\begin{xltabular}{\linewidth}{L{1.3cm}L{4.6cm}Y}
  \toprule
  \intestazione \tb{TC ID} & \tb{Descrizione Scenario} & \tb{Oracolo / Risultato Atteso} \\
  \midrule
  \endhead
  E2E-01 & \textbf{Ritorno dopo il login.} Un anonimo apre una pagina
  riservata. & Redirect al login; dopo l'accesso torna alla pagina
  richiesta. \\
  E2E-02 & \textbf{Password errata.} & Nessuna sessione; messaggio
  \enquote{Email o password non corretti. Riprova.} \\
  E2E-03 & \textbf{Studente nell'area staff.} & La pagina amministrativa non
  si apre e le API dello staff rispondono 403. \\
  E2E-04 & \textbf{Cookie inventato.} & Le pagine riservate portano al login,
  le API rispondono 401. \\
  E2E-05 & \textbf{Prenotazione guidata.} Data, mezza giornata, sala, posto. &
  Conferma della prenotazione, ritrovata tra \enquote{Le mie
  prenotazioni}. \\
  E2E-06 & \textbf{Posto occupato.} Un secondo studente apre la stessa
  fascia. & Il posto appare occupato e viene proposta la lista d'attesa. \\
  E2E-07 & \textbf{Ingresso in coda.} & Posizione 1, riletta dal server. \\
  E2E-08 & \textbf{Annullamento dello staff.} & L'operatore vede il nome
  dello studente promosso dalla coda. \\
  E2E-09 & \textbf{Notifica al promosso.} & Notifica \enquote{Posto assegnato
  dalla lista d'attesa} e prenotazione a suo nome. \\
  E2E-10 & \textbf{Conflitto all'ultimo istante.} Un secondo browser prende il
  posto durante la conferma. & Messaggio \enquote{Posto non più disponibile}
  con accesso alla lista d'attesa. \\
  E2E-11 & \textbf{Cancellazione del titolare.} & Chi è primo in coda ottiene
  il posto. \\
  E2E-12 & \textbf{Ingresso e uscita dalla coda.} Dalla scheda
  \enquote{Lista d'attesa}. & La scheda mostra la posizione 1 e il conteggio
  (1); dopo l'uscita compare \enquote{Non sei in nessuna lista d'attesa} e il
  conteggio torna a (0). \\
  \bottomrule
  \caption{Casi di Test E2E}\label{tab:e2e}\\
\end{xltabular}}

\end{document}
